\documentclass[12pt,a4paper]{article}
\usepackage[utf8]{inputenc}
\usepackage[french]{babel}
\usepackage{graphicx}
\usepackage{geometry}
\usepackage{listings}
\usepackage{xcolor}
\usepackage{hyperref}
\usepackage{fancyhdr}
\usepackage{titlesec}

\geometry{margin=2.5cm}

% Configuration des listings pour Java
\lstset{
    language=Java,
    basicstyle=\ttfamily\small,
    keywordstyle=\color{blue}\bfseries,
    commentstyle=\color{green!60!black}\itshape,
    stringstyle=\color{red},
    numbers=left,
    numberstyle=\tiny\color{gray},
    stepnumber=1,
    numbersep=8pt,
    backgroundcolor=\color{gray!10},
    frame=single,
    rulecolor=\color{black!30},
    breaklines=true,
    breakatwhitespace=false,
    tabsize=4,
    showstringspaces=false,
    captionpos=b
}

% En-tête et pied de page
\pagestyle{fancy}
\fancyhf{}
\rhead{Projet SMA - Contrôle Aérien}
\lhead{Master Sécurité IT \& Big Data}
\rfoot{Page \thepage}

\begin{document}

% Page de garde
\begin{titlepage}
    \centering
    \vspace*{2cm}

    {\LARGE\bfseries Université [Nom de l'Université]\par}
    \vspace{0.5cm}
    {\Large Master Sécurité IT \& Big Data\par}
    \vspace{2cm}

    {\Huge\bfseries Projet 8\par}
    \vspace{0.5cm}
    {\LARGE Application SMA pour le Contrôle Aérien\par}
    \vspace{3cm}

    {\Large\textbf{Module :} Systèmes Multi-Agents\par}
    \vspace{1.5cm}

    {\Large\textbf{Réalisé par :}\par}
    \vspace{0.3cm}
    {\large
        Ait Ezzaouite Mohamed\\
        Anas Ahrarche\\
        Nadir Mounim\par
    }
    \vspace{1.5cm}

    {\Large\textbf{Encadré par :}\par}
    \vspace{0.3cm}
    {\large Pr. El Mokhtar EN-NAIMI\par}

    \vfill

    {\large Année Universitaire 2024-2025\par}
\end{titlepage}

% Table des matières
\tableofcontents
\newpage

% Introduction
\section{Introduction}

La congestion du trafic aérien est un problème mondial qui coûte de plus en plus cher. Pour résoudre ce problème, il existe deux méthodes principales : la \textbf{gestion de la demande} et l'\textbf{amélioration de capacité}.

La gestion de la demande comprend l'allocation de créneaux horaires (slots) et la réduction de la surcharge pendant les heures de pointe. L'amélioration de la capacité peut être réalisée avec la construction de nouvelles pistes dans l'aéroport. Le coût élevé de la deuxième solution fait accroître l'intérêt de la méthode de gestion de la demande.

Le gestionnaire du trafic aérien modifie la demande en contrôlant les départs, de sorte que la capacité ne dépasse pas un seuil défini. La capacité s'incrémente en fournissant des séquences efficaces d'arrivées et départs d'avion pour minimiser leur temps d'inter arrivée et donc éliminer les temps idle. Cette tâche est essentiellement faite par le contrôleur de trafic aérien appelé \textbf{directeur de flux}.

Dans ce rapport, nous présentons notre système multi-agents (SMA) développé avec la plateforme JADE pour gérer efficacement le trafic aérien en allouant des créneaux horaires optimaux.

\section{Problématique}

Le principal défi dans la gestion du trafic aérien est d'optimiser l'utilisation des pistes d'atterrissage et de décollage tout en :
\begin{itemize}
    \item Respectant les priorités des vols (urgences, carburant faible, vols normaux)
    \item Minimisant les temps d'attente
    \item Évitant les conflits de créneaux horaires
    \item Gérant efficacement la capacité limitée des pistes
\end{itemize}

\section{Conception du Système Multi-Agents}

\subsection{Méthodologie Utilisée}

Notre conception s'appuie sur les principes des méthodologies AUML (Agent UML) et GAIA pour la modélisation des agents et de leurs interactions.

\subsection{Diagramme de Cas d'Utilisation}

Le diagramme de cas d'utilisation (Figure \ref{fig:usecase}) présente les principales fonctionnalités du système du point de vue des différents agents.

\begin{figure}[h]
    \centering
    \includegraphics[width=0.9\textwidth]{use-case.png}
    \caption{Diagramme de cas d'utilisation du système}
    \label{fig:usecase}
\end{figure}

Les acteurs principaux sont :
\begin{itemize}
    \item \textbf{FlightAgent} : Demande un créneau pour atterrir ou décoller
    \item \textbf{ControllerAgent} : Gère la file d'attente, trie par priorité et attribue les créneaux
    \item \textbf{RunwayAgent} : Vérifie la disponibilité des créneaux et confirme/rejette les propositions
\end{itemize}

\subsection{Diagramme de Classes}

Le diagramme de classes (Figure \ref{fig:classdiagram}) montre la structure statique de notre système multi-agents.

\begin{figure}[h]
    \centering
    \includegraphics[width=0.7\textwidth]{class-diagram.png}
    \caption{Diagramme de classes du système}
    \label{fig:classdiagram}
\end{figure}

\textbf{Description des classes principales :}

\begin{itemize}
    \item \textbf{FlightAgent} : Agent représentant un vol avec ses attributs (identifiant, priorité, heure demandée) et ses méthodes (demanderCreneau, recevoirAttribution)

    \item \textbf{ControllerAgent} : Agent contrôleur qui maintient une file d'attente de requêtes et implémente les méthodes de tri par priorité, attribution de créneaux et notification des vols

    \item \textbf{RunwayAgent} : Agent gérant le planning des pistes avec les méthodes de vérification de disponibilité et mise à jour du planning
\end{itemize}

\subsection{Diagramme de Séquence}

Le diagramme de séquence (Figure \ref{fig:sequencediagram}) illustre le flux d'interactions entre les agents lors d'une demande de créneau.

\begin{figure}[h]
    \centering
    \includegraphics[width=0.9\textwidth]{sequence-diagramm.png}
    \caption{Diagramme de séquence - Attribution de créneau}
    \label{fig:sequencediagram}
\end{figure}

\textbf{Scénario d'attribution de créneau :}
\begin{enumerate}
    \item Le \texttt{FlightAgent} envoie une demande de créneau au \texttt{ControllerAgent}
    \item Le \texttt{ControllerAgent} propose un créneau au \texttt{RunwayAgent}
    \item Le \texttt{RunwayAgent} accepte ou rejette la proposition
    \item Le \texttt{ControllerAgent} attribue le créneau et notifie le \texttt{FlightAgent}
    \item Le \texttt{FlightAgent} accuse réception
\end{enumerate}

\subsection{Diagramme d'Interaction}

Le diagramme d'interaction (Figure \ref{fig:interactiondiagram}) montre les échanges de messages entre les agents selon le protocole FIPA.

\begin{figure}[h]
    \centering
    \includegraphics[width=0.9\textwidth]{interaction-diagramm.png}
    \caption{Diagramme d'interaction entre agents}
    \label{fig:interactiondiagram}
\end{figure}

Les performatives FIPA utilisées sont :
\begin{itemize}
    \item \texttt{DEMANDE\_CRENEAU} : Requête initiale du vol
    \item \texttt{PROPOSER\_CRENEAU} : Proposition du contrôleur vers la piste
    \item \texttt{ACCEPTER / REJETER} : Réponse de la piste
    \item \texttt{CRENEAU\_PROPOSE / REFUS} : Notification au vol
    \item \texttt{ACCUSE\_RECEPTION} : Confirmation finale
\end{itemize}

\section{Architecture du Système}

\subsection{Rôles des Agents}

\subsubsection{FlightAgent (Agent Vol)}
\textbf{Rôle :} Représente un vol individuel (atterrissage ou décollage)

\textbf{Responsabilités :}
\begin{itemize}
    \item Envoyer une demande de créneau avec ses paramètres (ID, opération, heure souhaitée, priorité)
    \item Recevoir et confirmer l'attribution du créneau
\end{itemize}

\textbf{Comportements :}
\begin{itemize}
    \item OneShotBehaviour : Envoi de la requête initiale
    \item CyclicBehaviour : Écoute des réponses du contrôleur
\end{itemize}

\subsubsection{ControllerAgent (Directeur de Flux)}
\textbf{Rôle :} Gestionnaire central du trafic aérien

\textbf{Responsabilités :}
\begin{itemize}
    \item Maintenir une file d'attente des demandes de vols
    \item Trier les demandes par priorité puis par heure demandée
    \item Proposer des créneaux au RunwayAgent
    \item Gérer les rejets en reportant les créneaux (+5 minutes)
    \item Notifier les vols des attributions
\end{itemize}

\textbf{Algorithme de gestion :}
\begin{enumerate}
    \item Réception des nouvelles demandes
    \item Tri de la file : priorité (1 > 2 > 3) puis heure
    \item Arrondi de l'heure au multiple de 5 supérieur
    \item Négociation avec RunwayAgent
    \item Attribution ou report du créneau
\end{enumerate}

\subsubsection{RunwayAgent (Agent Piste)}
\textbf{Rôle :} Gestionnaire de la capacité de la piste

\textbf{Responsabilités :}
\begin{itemize}
    \item Maintenir le planning des créneaux occupés
    \item Vérifier la disponibilité des créneaux proposés
    \item Accepter ou rejeter les propositions
    \item Afficher l'état du planning
\end{itemize}

\subsection{Protocole de Communication}

Le système utilise le protocole FIPA-Request avec les performatives ACL suivantes :
\begin{itemize}
    \item \texttt{REQUEST} : Demande de créneau (Flight → Controller)
    \item \texttt{PROPOSE} : Proposition de créneau (Controller → Runway)
    \item \texttt{ACCEPT\_PROPOSAL} : Créneau disponible (Runway → Controller)
    \item \texttt{REJECT\_PROPOSAL} : Créneau occupé (Runway → Controller)
    \item \texttt{INFORM} : Attribution confirmée (Controller → Flight)
\end{itemize}

\section{Implémentation}

\subsection{Plateforme Utilisée}

Nous avons choisi \textbf{JADE} (Java Agent DEvelopment Framework) comme plateforme de développement pour les raisons suivantes :
\begin{itemize}
    \item Conformité aux standards FIPA
    \item Bibliothèques riches pour la communication inter-agents
    \item Facilité de déploiement et de débogage
    \item Large communauté et documentation
\end{itemize}

\subsection{Code Source}

\subsubsection{Main.java}

Point d'entrée du système qui initialise le conteneur JADE et crée les agents.

\begin{lstlisting}[caption=Main.java]
import jade.core.Profile;
import jade.core.ProfileImpl;
import jade.core.Runtime;
import jade.wrapper.AgentContainer;

public class Main {
    public static void main(String[] args) throws Exception {
        Runtime rt = Runtime.instance();
        Profile p = new ProfileImpl();
        AgentContainer mc = rt.createMainContainer(p);

        // Runway
        mc.createNewAgent("runway", "RunwayAgent", null).start();

        // Controller
        mc.createNewAgent("controller", "ControllerAgent", null).start();

        // Flights (hardcoded)
        mc.createNewAgent("F1", "FlightAgent", new Object[]{"F1","LAND","10","1"}).start();
        mc.createNewAgent("F2", "FlightAgent", new Object[]{"F2","TAKEOFF","12","3"}).start();
        mc.createNewAgent("F3", "FlightAgent", new Object[]{"F3","LAND","11","2"}).start();
        mc.createNewAgent("F4", "FlightAgent", new Object[]{"F4","LAND","10","3"}).start();
        mc.createNewAgent("F5", "FlightAgent", new Object[]{"F5","TAKEOFF","10","1"}).start();
    }
}
\end{lstlisting}

\subsubsection{FlightAgent.java}

Agent représentant un vol individuel.

\begin{lstlisting}[caption=FlightAgent.java]
import jade.core.Agent;
import jade.core.behaviours.OneShotBehaviour;
import jade.core.behaviours.CyclicBehaviour;
import jade.lang.acl.ACLMessage;
import jade.core.AID;

public class FlightAgent extends Agent {
    protected void setup() {
        Object[] args = getArguments(); // id, op, time, priority
        String flightId = (String) args[0];
        String op = (String) args[1];
        int time = Integer.parseInt(args[2].toString());
        int priority = Integer.parseInt(args[3].toString());

        addBehaviour(new OneShotBehaviour() {
            public void action() {
                ACLMessage req = new ACLMessage(ACLMessage.REQUEST);
                req.addReceiver(new AID("controller", AID.ISLOCALNAME));
                req.setContent(flightId + "," + op + "," + time + "," + priority);
                send(req);
            }
        });

        addBehaviour(new CyclicBehaviour() {
            public void action() {
                ACLMessage msg = receive();
                if (msg != null && msg.getPerformative() == ACLMessage.INFORM) {
                    System.out.println(getLocalName() + " received: " + msg.getContent());
                } else {
                    block();
                }
            }
        });
    }
}
\end{lstlisting}

\subsubsection{ControllerAgent.java}

Agent contrôleur qui gère la file d'attente et l'attribution des créneaux.

\begin{lstlisting}[caption=ControllerAgent.java]
import jade.core.Agent;
import jade.core.behaviours.CyclicBehaviour;
import jade.lang.acl.ACLMessage;
import jade.core.AID;
import jade.lang.acl.MessageTemplate;

import java.util.*;

public class ControllerAgent extends Agent {
    private List<FlightRequest> queue = new ArrayList<>();
    private AID runwayAgent;

    protected void setup() {
        runwayAgent = new AID("runway", AID.ISLOCALNAME);

        addBehaviour(new CyclicBehaviour() {
            public void action() {
                // Receive new requests
                ACLMessage msg = receive();
                if (msg != null && msg.getPerformative() == ACLMessage.REQUEST) {
                    String[] p = msg.getContent().split(",");
                    FlightRequest fr = new FlightRequest(
                            p[0], p[1], Integer.parseInt(p[2]), Integer.parseInt(p[3])
                    );
                    queue.add(fr);
                }

                if (!queue.isEmpty()) {
                    // Sort by priority then requestedTime
                    queue.sort(Comparator
                            .comparingInt((FlightRequest f) -> f.priority)
                            .thenComparingInt(f -> f.requestedTime));

                    FlightRequest top = queue.get(0);
                    int slotTime = roundUpTo5(top.requestedTime);

                    ACLMessage propose = new ACLMessage(ACLMessage.PROPOSE);
                    propose.addReceiver(runwayAgent);
                    propose.setContent(top.flightId + "," + slotTime);
                    send(propose);

                    // Wait only for runway response
                    MessageTemplate mt = MessageTemplate.MatchSender(runwayAgent);
                    ACLMessage reply = blockingReceive(mt);

                    if (reply.getPerformative() == ACLMessage.ACCEPT_PROPOSAL) {
                        queue.remove(top);

                        // Notify flight
                        ACLMessage inform = new ACLMessage(ACLMessage.INFORM);
                        inform.addReceiver(new AID(top.flightId, AID.ISLOCALNAME));
                        inform.setContent("SLOT_ASSIGNED," + slotTime);
                        send(inform);

                        System.out.println("Assigned " + top.flightId +
                                " -> slot " + slotTime + " (priority " + top.priority + ")");
                    } else {
                        // If rejected, push request to next slot (+5)
                        top.requestedTime += 5;
                    }
                } else {
                    block(200);
                }
            }
        });
    }

    private int roundUpTo5(int t) {
        return ((t + 4) / 5) * 5;
    }
}
\end{lstlisting}

\subsubsection{RunwayAgent.java}

Agent gérant la capacité et le planning de la piste.

\begin{lstlisting}[caption=RunwayAgent.java]
import jade.core.Agent;
import jade.core.behaviours.CyclicBehaviour;
import jade.lang.acl.ACLMessage;

import java.util.HashMap;
import java.util.Map;

public class RunwayAgent extends Agent {
    private Map<Integer, String> schedule = new HashMap<>(); // slotTime -> flightId

    protected void setup() {
        addBehaviour(new CyclicBehaviour() {
            public void action() {
                ACLMessage msg = receive();
                if (msg != null && msg.getPerformative() == ACLMessage.PROPOSE) {
                    String[] parts = msg.getContent().split(",");
                    String flightId = parts[0];
                    int slotTime = Integer.parseInt(parts[1]);

                    ACLMessage reply = msg.createReply();
                    if (!schedule.containsKey(slotTime)) {
                        schedule.put(slotTime, flightId);
                        reply.setPerformative(ACLMessage.ACCEPT_PROPOSAL);
                        reply.setContent("ACCEPTED");

                        // Print current schedule
                        System.out.println("Runway schedule: " + schedule);
                    } else {
                        reply.setPerformative(ACLMessage.REJECT_PROPOSAL);
                        reply.setContent("REJECTED");
                    }
                    send(reply);
                } else {
                    block();
                }
            }
        });
    }
}
\end{lstlisting}

\subsubsection{FlightRequest.java}

Classe de données représentant une demande de vol.

\begin{lstlisting}[caption=FlightRequest.java]
public class FlightRequest {
    public String flightId;
    public String operation; // "LAND" or "TAKEOFF"
    public int requestedTime; // minutes since 0
    public int priority; // 1=Emergency, 2=LowFuel, 3=Normal

    public FlightRequest(String id, String op, int time, int priority) {
        this.flightId = id;
        this.operation = op;
        this.requestedTime = time;
        this.priority = priority;
    }
}
\end{lstlisting}

\subsubsection{Slot.java}

Classe représentant un créneau horaire.

\begin{lstlisting}[caption=Slot.java]
public class Slot {
    public int time; // slot time in minutes
    public String flightId;

    public Slot(int time, String flightId) {
        this.time = time;
        this.flightId = flightId;
    }
}
\end{lstlisting}

\section{Fonctionnement du Système}

\subsection{Algorithme de Gestion des Priorités}

Le système utilise un algorithme de tri multi-critères :

\begin{enumerate}
    \item \textbf{Critère principal :} Priorité du vol
    \begin{itemize}
        \item Priorité 1 : Urgence (problème technique, urgence médicale)
        \item Priorité 2 : Carburant faible
        \item Priorité 3 : Vol normal
    \end{itemize}

    \item \textbf{Critère secondaire :} Heure de demande (ordre chronologique)

    \item \textbf{Arrondi des créneaux :} Tous les créneaux sont arrondis au multiple de 5 minutes supérieur pour standardiser le planning
\end{enumerate}

\subsection{Gestion des Conflits}

Lorsqu'un créneau proposé est déjà occupé :
\begin{enumerate}
    \item Le RunwayAgent rejette la proposition
    \item Le ControllerAgent reporte la demande de 5 minutes
    \item Le vol reste dans la file d'attente pour une nouvelle tentative
    \item Le processus se répète jusqu'à trouver un créneau disponible
\end{enumerate}

\subsection{Exemple d'Exécution}

Pour le scénario défini dans \texttt{Main.java} :

\begin{verbatim}
F1: LAND, temps=10, priorité=1 (urgence)
F2: TAKEOFF, temps=12, priorité=3 (normal)
F3: LAND, temps=11, priorité=2 (carburant faible)
F4: LAND, temps=10, priorité=3 (normal)
F5: TAKEOFF, temps=10, priorité=1 (urgence)
\end{verbatim}

\textbf{Ordre d'attribution :}
\begin{enumerate}
    \item F1 → créneau 10 (priorité 1, demandé à t=10)
    \item F5 → créneau 15 (priorité 1, demandé à t=10, conflit donc reporté)
    \item F3 → créneau 15 (priorité 2, demandé à t=11)
    \item F4 → créneau 20 (priorité 3, demandé à t=10)
    \item F2 → créneau 15 (priorité 3, demandé à t=12)
\end{enumerate}

\section{Tests et Validation}

\subsection{Scénarios de Test}

Nous avons testé plusieurs scénarios :

\begin{enumerate}
    \item \textbf{Test de priorité :} Vérification que les urgences sont traitées en premier
    \item \textbf{Test de conflit :} Gestion correcte des créneaux déjà occupés
    \item \textbf{Test de charge :} Comportement avec un grand nombre de vols
    \item \textbf{Test de communication :} Bon échange de messages FIPA
\end{enumerate}

\subsection{Résultats}

Le système fonctionne correctement et respecte les contraintes suivantes :
\begin{itemize}
    \item[\checkmark] Attribution basée sur les priorités
    \item[\checkmark] Pas de double allocation de créneau
    \item[\checkmark] Gestion automatique des conflits
    \item[\checkmark] Communication asynchrone efficace
\end{itemize}

\section{Avantages et Limites}

\subsection{Avantages}

\begin{itemize}
    \item \textbf{Décentralisation :} Chaque agent gère sa propre logique
    \item \textbf{Scalabilité :} Facilement extensible à plusieurs pistes
    \item \textbf{Robustesse :} Gestion automatique des conflits
    \item \textbf{Standards :} Conformité FIPA et JADE
\end{itemize}

\subsection{Limites et Améliorations Possibles}

\begin{itemize}
    \item \textbf{Piste unique :} Actuellement limité à une seule piste
    \item \textbf{Vols hardcodés :} Pas d'interface pour ajouter dynamiquement des vols
    \item \textbf{Pas de GUI :} Interface en ligne de commande uniquement
    \item \textbf{Algorithme simple :} Possibilité d'optimiser davantage l'allocation
\end{itemize}

\textbf{Améliorations futures :}
\begin{itemize}
    \item Ajouter une interface graphique (JADE GUI ou web)
    \item Gérer plusieurs pistes simultanément
    \item Implémenter des algorithmes d'optimisation avancés
    \item Ajouter des statistiques de performance
    \item Intégrer des données météorologiques
\end{itemize}

\section{Conclusion}

Ce projet nous a permis de développer un système multi-agents fonctionnel pour la gestion du trafic aérien en utilisant la plateforme JADE. Nous avons appliqué les concepts des systèmes multi-agents (autonomie, communication, négociation) pour résoudre un problème complexe de gestion de ressources partagées.

Le système démontre l'efficacité de l'approche multi-agents pour :
\begin{itemize}
    \item La coordination décentralisée
    \item La gestion de priorités multiples
    \item La résolution de conflits de ressources
    \item La communication asynchrone entre entités autonomes
\end{itemize}

Les méthodologies AUML et GAIA nous ont fourni un cadre structuré pour la conception, tandis que JADE a facilité l'implémentation avec ses bibliothèques conformes aux standards FIPA.

Ce travail constitue une base solide qui pourrait être étendue vers un système de gestion du trafic aérien plus complet et réaliste.

\section*{Références}

\begin{itemize}
    \item FIPA Agent Communication Language Specifications
    \item JADE Programming Tutorial - Fabio Bellifemine et al.
    \item Méthodologie GAIA pour les Systèmes Multi-Agents
    \item AUML (Agent UML) - Extension d'UML pour les agents
    \item Documentation officielle JADE : \url{https://jade.tilab.com/}
\end{itemize}

\end{document}
